{% extends "members/tex_base.tex" %}
{% load static common tex_extras %}

{% block title %}Abrechnungs- und Zuschussbeleg\\[2mm]Sektionsveranstaltung{% endblock %}

{% block content %}

{% if excursion %}
\noindent\textbf{\large Ausfahrt}

% DESCRIPTION TABLE
\begin{table}[H]
    \begin{tabular}{ll}
        Aktivität: & {{ excursion.name|esc_all }} \\
        Ordnungsnummer & {{ excursion.code|esc_all }} \\
        Ort / Stützpunkt: & {{ excursion.place|esc_all }} \\
        Zeitraum: & {{ excursion.duration|esc_all}} Tage ({{ excursion.time_period_str|esc_all }}) \\
        Teilnehmer*innen: & {{ excursion.participant_count }} der Gruppe(n) {{ excursion.groups_str|esc_all }} \\
        Betreuer*innen: & {{excursion.staff_count|esc_all }} ({{ excursion.staff_str|esc_all }}) \\
        Art der Tour: & {% checked_if_true 'Gemeinschaftstour' excursion.get_tour_type %}
            {% checked_if_true 'Führungstour' excursion.get_tour_type %}
            {% checked_if_true 'Ausbildung' excursion.get_tour_type %} \\
        Anreise: & {% checked_if_true 'ÖPNV' excursion.get_tour_approach %}
            {% checked_if_true 'Muskelkraft' excursion.get_tour_approach %}
            {% checked_if_true 'Fahrgemeinschaften' excursion.get_tour_approach %}
    \end{tabular}
\end{table}

\noindent\textbf{\large Zuschüsse und Aufwandsentschädigung}
{% if excursion.approved_staff_count > 0 %}

\noindent Gemäß Beschluss des Jugendausschusses gelten folgende Sätze für Zuschüsse pro genehmigter Jugendleiter*in:

\begin{table}[H]
    \centering
    \begin{tabularx}{.97\textwidth}{Xllr}
        \toprule
        \textbf{Posten} & \textbf{Einzelsatz} & \textbf{Anzahl} & \textbf{Gesamtbetrag pro JL} \\
        \midrule
        Zuschuss Übernachtung & {{ statement.price_per_night }} € / Nacht  & {{ statement.nights }} Nächte & {{ statement.nights_per_yl }} € \\
        Zuschuss Anreise & {{statement.euro_per_km}} € / km ({{ statement.means_of_transport }}) & {{ statement.kilometers_traveled }} km & {{ statement.transportation_per_yl }} € \\
        Aufwandsentschädigung & {{ statement.allowance_per_day }},00 € / Tag  & {{ statement.duration }} Tage & {{ statement.allowance_per_yl }} € \\
        \midrule
        \textbf{Summe}& & & \textbf{ {{ statement.total_per_yl }} }€\\
        \bottomrule
    \end{tabularx}
\end{table}

\noindent Gemäß JDAV-Betreuungsschlüssel können bei {{ excursion.participant_count }} Teilnehmer*innen
bis zu {{ excursion.approved_staff_count }} Jugendleiter*innen {% if excursion.approved_extra_youth_leader_count %}
(davon {{ excursion.approved_extra_youth_leader_count }} durch das Jugendreferat zusätzlich genehmigt){% endif %} bezuschusst werden.
Zuschüsse und Aufwandsentschädigung werden wie folgt abgerufen:
% One item per line the section's rules resolved, in the order they resolved
% it: the label, the arithmetic and the recipient are the rules' own words.
% Nothing in this document knows what an allowance or an LJP contribution is.
\begin{itemize}
    {% for line in statement.payout %}
    \item {% if line.negative %}abzüglich {% endif %}{{ line.label|esc_all }} in Höhe von {{ line.amount_abs }} €{% if line.recipient %} {% if line.negative %}von{% else %}an{% endif %} {{ line.recipient|esc_all }}{% endif %}{% if line.explanation %} ({{ line.explanation|esc_all }}){% endif %}
    {% empty %}
    \item Es werden keine Zuschüsse und keine Aufwandsentschädigung in Anspruch genommen.
    {% endfor %}
\end{itemize}
{% else %}
\noindent Für die vorliegende Ausfahrt sind keine Jugendleiter*innen anspruchsberechtigt für Zuschüsse oder Aufwandsentschädigung.

{% endif %}

{% else %}
\vspace{110pt}
{% endif %}


\vspace{12pt}

\noindent\textbf{\large Ausgabenübersicht}
\nopagebreak
\begin{table}[H]
    \centering
    \begin{tabularx}{.97\textwidth}{lXlr}
        \toprule
        \textbf{Titel} & \textbf{Beschreibung} & \textbf{Auszahlung an} & \textbf{Betrag} \\
        \midrule

{% if statement.bills_covered %}
    {% for bill in statement.bills_covered %}
        {{ forloop.counter }}. {{ bill.short_description}} & {{ bill.explanation}} & {{ bill.paid_by.name|esc_all }} &  {{ bill.amount }} € \\
        {% endfor %}
        \midrule
        \multicolumn{3}{l}{\textbf{Summe übernommene Ausgaben}} & \textbf{ {{ statement.total_bills }} }€\\
{% endif %}
{% if statement.payout %}
        \midrule
        % Every line the section's rules resolved, in the order they resolved
        % it. Nothing here knows what an allowance or an LJP contribution is.
        {% for line in statement.payout %}
        \multicolumn{2}{l}{ {{ line.label|esc_all }}{% if line.explanation %} ({{ line.explanation|esc_all }}){% endif %} } & {{ line.recipient|esc_all }} & {{ line.amount }} €\\
        {% endfor %}
        \midrule
        \multicolumn{3}{l}{\textbf{Summe Zuschüsse und Entschädigungen}} & \textbf{ {{ statement.payout_total }} }€\\
{% endif %}
{% if statement.payout or statement.bills_covered %}
        \midrule
        \textbf{Gesamtsumme}& & & \textbf{ {{ statement.total }} }€\\
{% endif %}
        \bottomrule
    \end{tabularx}
\end{table}

% The payment orders, and only where this document *is* the instruction: in a
% section whose Geschäftsstelle makes the transfers, nobody who makes them ever
% sees the review screen, so the accounts, the references and the amounts have
% to be on the paper. Where the reviewer transfers the money from that screen,
% `handoff` is false and this section does not exist -- the document is then a
% receipt for transfers that have already been made.
{% if statement.handoff %}
\vspace{12pt}

\noindent\textbf{\large Überweisungsaufträge}

\noindent Die folgenden Überweisungen sind von der Geschäftsstelle auszuführen.
{% if statement.orders %}

\begin{table}[H]
    \centering
    \begin{tabularx}{.97\textwidth}{lXr}
        \toprule
        \textbf{Empfänger*in} & \textbf{IBAN und Verwendungszweck} & \textbf{Betrag} \\
        \midrule
        % One row per booking, which is what `reduce_transactions` has already
        % bundled -- so this is the set of transfers to make and not a list to
        % be combined again by whoever reads it. `separate` marks the rows that
        % must *not* be combined, with the rules' own name for the group as the
        % reason.
        {% for order in statement.orders %}
        {{ order.recipient|esc_all }} & {% if order.iban %}\texttt{ {{ order.iban|esc_all }} }{% else %}\textbf{keine IBAN hinterlegt}{% endif %} \newline {{ order.reference|esc_all }}{% if order.separate %} \newline \emph{getrennt zu überweisen{% if order.settlement_group %} (Abrechnungsgruppe {{ order.settlement_group|esc_all }}){% endif %}}{% endif %} & {{ order.amount }} €\\
        {% endfor %}
        \midrule
        \multicolumn{2}{l}{\textbf{Summe Überweisungen}} & \textbf{ {{ statement.orders_total }} }€\\
        \bottomrule
    \end{tabularx}
\end{table}

\noindent Zeilen mit dem Hinweis \emph{getrennt zu überweisen} dürfen nicht zu einer
Überweisung zusammengefasst werden, auch wenn sie an dieselbe Person gehen: sie
gehören zu unterschiedlichen Abrechnungsgruppen, etwa weil eine davon steuerlich
getrennt nachweisbar sein muss. Alle übrigen Zeilen sind bereits zusammengefasst.
{% else %}

\noindent Für diese Abrechnung sind keine Überweisungen auszuführen.
{% endif %}
{% endif %}

\noindent Dieser Beleg wird automatisch erstellt und daher nicht unterschrieben.

{% endblock %}
